\begin{cuebabstract}
本科毕业论文需要同时满足内容论证与形式表达的要求。为了帮助使用者理解论文源文件的组织方式，本示例以排版流程为对象，展示基本信息、摘要、三级标题、数学公式、图表、文献脚注、参考文献、附录和致谢的相互关系。示例采用分文件组织方法，将正文内容与公共格式配置分开；通过标签和交叉引用维护对象之间的联系，并用显式语言字段及作者拼音排序键组织文献。本文件中的数值均为人工构造的演示数据，只用于观察表格、公式和分页效果，不代表任何调查、实验或学校统计结果。示例进一步给出从内容检查、编译检查到页面核对的工作顺序，说明编译成功与符合正式提交要求之间的区别。当前版本依据所附二〇二四届手册建立通用排版原型，保留对封面底稿、字体、装订线和编号歧义的核对事项。全文为缩短的教学示例，不能替代满足学校字数要求的毕业论文，也不构成正式格式认证。
\cuebkeywords{毕业论文；中文排版；文献管理；格式核对}
\end{cuebabstract}

\begin{cuebabstract*}
An undergraduate thesis requires both a reasoned argument and a consistent presentation. This example introduces a writing workflow through a small document containing metadata, abstracts, three levels of headings, equations, figures, tables, citation footnotes, references, an appendix, and acknowledgements. Content is stored separately from common formatting settings. Labels and cross references connect document objects, while explicit language fields and author sorting keys support bibliography ordering. Every numerical value in the example is synthetic and serves only to demonstrate typesetting; none represents a survey, experiment, or institutional statistic. The example also outlines a review sequence covering source content, compilation, and page inspection. Successful compilation alone does not establish compliance with submission requirements. The current preview follows the supplied handbook for the 2024 graduating class and records open questions about the approved cover, fonts, binding allowance, and numbering rules. This shortened teaching document is neither a complete graduation thesis nor an official certification of its format.
\cuebkeywords{undergraduate thesis; Chinese typesetting; reference management; format review}
\end{cuebabstract*}
